\documentclass[11pt,a4paper]{letter}

\usepackage[T1]{fontenc}
\usepackage[utf8]{inputenc}
\usepackage{lmodern}
\usepackage[margin=1in]{geometry}
\usepackage{hyperref}
\hypersetup{hidelinks}

\signature{Darcio Genicolo-Martins \\ Paulo Furquim de Azevedo \\ Insper Institute of Education and Research \\ S\~ao Paulo, Brazil}
\address{Darcio Genicolo-Martins \\ Paulo Furquim de Azevedo \\ Insper Institute of Education and Research \\ Rua Quat\'a, 300 \\ S\~ao Paulo, SP 04546-042 \\ Brazil \\ \href{mailto:darciogm1@insper.edu.br}{darciogm1@insper.edu.br}}
\date{April 17, 2026}

\begin{document}

\begin{letter}{The Editors \\ Journal of Public Economics \\ Elsevier}

\opening{Dear Editors,}

We are pleased to submit the manuscript \emph{``Bitter Pills: Judicial Enforcement and the Cost of Public Procurement in Brazil''} for consideration as a \textbf{Short Paper} in the \emph{Journal of Public Economics}.

\medskip\noindent\textbf{Contribution.} The paper provides the first estimates of the fiscal cost of judicial enforcement in public procurement. Using bid-level administrative data on 479{,}330 pharmaceutical purchases by the S\~ao Paulo State Department of Health over 2009--2019, we show that court-mandated procurement raises per-unit prices by 5.4\%, reduces bidder participation by 5.4\%, and nonetheless closes tenders at higher rates---a pattern consistent with officials accepting worse terms to avoid personal sanctions.

\medskip\noindent\textbf{Identification.} Our institutional setting includes a quasi-unique feature: a parallel \emph{administrative-request} channel that shares all planning constraints with litigated procurement but carries no penalties for procurement failure. Comparing litigated and administrative urgent purchases within the same item---what we call the ``under the gun'' comparison---isolates a 23--30\% cost gap attributable to the sanction regime. We frame this as a descriptive decomposition rather than a structural causal parameter, discuss threats to identification explicitly, validate parallel trends using the honest-DiD imputation estimator of Borusyak-Jaravel-Spiess (2024), and falsify the premium on items never subject to litigation.

\medskip\noindent\textbf{Mechanism.} A supplier-fixed-effects decomposition splits the premium roughly in half between demand-side selection (officials accept worse matches under pressure) and within-firm markups (the same suppliers charge more when the government has no alternative), with distinct policy implications for each channel. Brazil's recent procurement reform (Lei~14.133/2021), which introduces framework agreements, could mitigate the quantity-fragmentation channel we document.

\medskip\noindent\textbf{JPubE fit.} The paper sits at the intersection of procurement design, judicial enforcement, and fiscal efficiency---themes central to the journal's agenda. We comply with the Short Paper track constraints: the body totals 5{,}370 words (well under the 6{,}000-word cap), uses 5 exhibits (4 tables and 1 figure), and is self-contained---the main conclusions do not depend on the Online Appendix.

\medskip\noindent\textbf{Supplementary material.} The Online Appendix (submitted separately) contains descriptive statistics, within-item balance tests, honest-DiD estimates, sensitivity to winsorization, heterogeneity by item class and market competition, and the hand-labeled validation of our regex-based classification.

\medskip\noindent\textbf{Declarations.} The work is original, has not been published elsewhere, and is not under consideration at any other journal. The authors report no conflicts of interest. Replication code and de-identified data files will be shared upon request, subject to the confidentiality restrictions imposed by the S\~ao Paulo State Department of Health.

We defer to your editorial judgment on referees but would gladly suggest potential reviewers if helpful. Thank you for considering our submission.

\closing{Sincerely,}

\end{letter}

\end{document}
